\documentclass[11pt]{article}
\usepackage[a4paper,margin=1in]{geometry}
\usepackage{amsmath,amssymb,amsfonts,bm}
\usepackage{booktabs,array,longtable}
\usepackage[hidelinks]{hyperref}
\usepackage{enumitem}
\usepackage{microtype}
\usepackage{mathtools}
\setlength{\parskip}{0.45em}
\setlength{\parindent}{0pt}

\title{Appendix: Method, Units, and Interpretation of the Revised SCWF Pipeline}
\author{}
\date{March 11, 2026}

\begin{document}
\maketitle

\section{Scope}

This appendix documents the \emph{implemented} method in the revised pipeline
contained in \path{three_D_funcs.py}. Its purpose is narrow. It states the actual
filter equations used by the code, the associated units, the statistical objects
that are estimated, and the limits of their interpretation.

The central point is the following. The fluctuating quantities in this pipeline are
\emph{wavelet coefficients}, not finite-difference increments. The code preserves
legacy names such as \texttt{dB}, \texttt{dV}, \texttt{dZp}, and \texttt{dZm} for
compatibility, but their mathematically correct interpretation is a scale-local
coefficient family.

\section{Wavelet-frame specification}

Let $\Delta t$ denote the cadence. For derivative order $M$ and voices-per-octave
$\nu$, the code defines the scale family
\begin{equation}
 s_j = s_0 \, 2^{(j-1)/\nu},
\end{equation}
with
\begin{equation}
 s_0 = \max\!\left(1, \frac{\sqrt{M}}{2\pi}\right)
\end{equation}
in sample units and therefore
\begin{equation}
 s_j^{(\mathrm{sec})} = s_j \, \Delta t.
\end{equation}

The parser for the wavelet token follows the code, not an external naming
convention. In particular, \texttt{mw8} means \emph{8 voices per octave}. Unless an
explicit order token is present, the derivative order defaults to $M=4$.

\section{Implemented low-pass and band-pass responses}

The low-pass background response is the Gaussian
\begin{equation}
 \widehat{L}_j(\omega) = \exp\!\left[-\frac{1}{2}(s_j^{(\mathrm{sec})}\omega)^2\right].
\end{equation}

The band-pass coefficient response is the even derivative-of-Gaussian
\begin{equation}
 \widehat{H}_j(\omega) = \frac{\sqrt{s_j^{(\mathrm{sec})}}\,(s_j^{(\mathrm{sec})}\omega)^M
 \exp\!\left[-\frac{1}{2}(s_j^{(\mathrm{sec})}\omega)^2\right]}{\|\psi_M\|_+},
\end{equation}
with one-sided mother norm
\begin{equation}
 \|\psi_M\|_+^2 = \int_0^{\infty} \eta^{2M} e^{-\eta^2} \, d\eta.
\end{equation}

The code uses real FFTs and therefore works with the positive-frequency branch of
the response. The low-pass and band-pass channels are \emph{not} a complementary
partition of unity. In general,
\begin{equation}
 x(t) \neq x_{l,j}(t) + W_j(t).
\end{equation}
The background and coefficient fields are paired multiscale diagnostics rather than
an exact additive decomposition.

\section{Implemented transforms and coefficient units}

For any scalar or vector field $X(t)$, the code evaluates
\begin{equation}
 X_{l,j}(t) = \mathcal{F}^{-1}\!\left[\widehat{L}_j(\omega)\,\widehat{X}(\omega)\right],
\end{equation}
and
\begin{equation}
 W_{X,j}(t) = \mathcal{F}^{-1}\!\left[\widehat{H}_j(\omega)\,\widehat{X}(\omega)\right].
\end{equation}
Because the responses are even in frequency and the padding is symmetric, the
interior estimator is centered and zero-phase.

If $X$ has physical units $[X]$, then
\begin{equation}
 [X_{l,j}] = [X], \qquad [W_{X,j}] = [X]\,\mathrm{s}^{1/2}.
\end{equation}
This $\mathrm{s}^{1/2}$ factor is the direct consequence of the implemented
$L^2$-type wavelet normalization.

\section{Lag calibration used by the code}

The peak-response frequency of the implemented band-pass filter is
\begin{equation}
 f_{\mathrm{peak},j} = \frac{\sqrt{M}}{2\pi s_j^{(\mathrm{sec})}},
\end{equation}
and the code defines the corresponding equivalent lag as
\begin{equation}
 \tau_j = \frac{1}{f_{\mathrm{peak},j}} = \frac{2\pi s_j^{(\mathrm{sec})}}{\sqrt{M}}.
\end{equation}

This is a calibration of the filter center. It is not an exact finite-difference
lag.

\section{Gap handling and cone of influence}

The FFT backend cannot tolerate NaNs. The code therefore interpolates gaps only for
the purpose of evaluating the transform, then applies explicit validity masks based
on the original finite-sample support.

For a support half-width $b_j$ in samples, a coefficient at index $t$ is retained
only if
\begin{equation}
 t \in [b_j, N-b_j)
\end{equation}
and the raw data contain no invalid sample inside that support neighborhood. In the
implementation,
\begin{equation}
 b_j = \left\lceil \sigma_{\mathrm{sup}} s_j \right\rceil,
\end{equation}
with $\sigma_{\mathrm{sup}}$ the configured support multiplier.

\section{Magnetic, velocity, and Elsasser coefficient families}

For the magnetic field $\mathbf{B}$, velocity $\mathbf{V}$, and density $N_p$, the
code computes scale-local backgrounds
\begin{equation}
 \mathbf{B}_{l,j}(t), \qquad \mathbf{V}_{l,j}(t), \qquad N_{l,j}(t),
\end{equation}
and coefficient fields
\begin{equation}
 \mathbf{W}_{B,j}(t), \qquad \mathbf{W}_{V,j}(t).
\end{equation}

The magnetic coefficient is converted to Alfv\'enic velocity units through
\begin{equation}
 \mathbf{W}_{B_v,j}(t) = C_A\!\left(N_{l,j}(t)\right) \, \mathbf{W}_{B,j}(t),
\end{equation}
with
\begin{equation}
 C_A(N) = \frac{10^{-15}}{\sqrt{\mu_0 m_p N}}.
\end{equation}
The code assumes $N$ is in cm$^{-3}$ and returns $C_A$ in km s$^{-1}$ per nT.
Therefore $\mathbf{W}_{B_v,j}$ has units km s$^{-1}$ s$^{1/2}$.

If spacecraft-velocity correction is requested, the same corrected velocity is used
consistently both in the Taylor mapping and in the velocity coefficient family used
for Elsasser construction.

With either the global or local polarity sign $s_B(t)$, the coefficient-based
Elsasser variables are
\begin{equation}
 \mathbf{W}_{Z^{\pm},j}(t) = \mathbf{W}_{V,j}(t) \pm s_B(t)\,\mathbf{W}_{B_v,j}(t).
\end{equation}
These are coefficient surrogates in Elsasser variables. They are not literal
finite-difference Elsasser increments.

\section{Local anisotropy basis and projected separations}

The code assigns a Taylor-mapped local separation vector
\begin{equation}
 \bm{\ell}_j(t) = \mathbf{V}_{l,j}(t)\,\tau_j.
\end{equation}

The local basis is then built from the local background field and the perpendicular
magnetic coefficient direction,
\begin{equation}
 \hat{\mathbf{e}}_{\ell} = \frac{\mathbf{B}_{l,j}}{|\mathbf{B}_{l,j}|},
\end{equation}
\begin{equation}
 \mathbf{W}_{B,\perp,j} = \mathbf{W}_{B,j} - \frac{\mathbf{W}_{B,j}\cdot \mathbf{B}_{l,j}}{|\mathbf{B}_{l,j}|^2}\,\mathbf{B}_{l,j},
\end{equation}
\begin{equation}
 \hat{\mathbf{e}}_{\xi} = \frac{\mathbf{W}_{B,\perp,j}}{|\mathbf{W}_{B,\perp,j}|},
\end{equation}
\begin{equation}
 \hat{\mathbf{e}}_{\lambda} = \frac{\hat{\mathbf{e}}_{\ell}\times\hat{\mathbf{e}}_{\xi}}{|\hat{\mathbf{e}}_{\ell}\times\hat{\mathbf{e}}_{\xi}|}.
\end{equation}

The projected separation magnitudes are
\begin{equation}
 \ell_{\ell,j} = \frac{|\bm{\ell}_j\cdot\hat{\mathbf{e}}_{\ell}|}{d_i}, \qquad
 \ell_{\xi,j} = \frac{|\bm{\ell}_j\cdot\hat{\mathbf{e}}_{\xi}|}{d_i}, \qquad
 \ell_{\lambda,j} = \frac{|\bm{\ell}_j\cdot\hat{\mathbf{e}}_{\lambda}|}{d_i},
\end{equation}
and
\begin{equation}
 \ell_{\mathrm{mag},j} = \frac{|\bm{\ell}_j|}{d_i}.
\end{equation}

The basis is undefined when $|\mathbf{B}_{l,j}|=0$ or $|\mathbf{W}_{B,\perp,j}|=0$.
The implementation masks those samples rather than assigning an artificial
orientation.

\section{What the angle buckets actually represent}

The angle thresholds define conditioned subsets such as
\begin{align}
 \texttt{ell\_par}:& \qquad \theta < \theta_{\parallel}, \\
 \texttt{Ell\_perp}:& \qquad \theta > \theta_{\xi}, \quad \phi < \phi_{\xi}, \\
 \texttt{ell\_perp}:& \qquad \theta > \theta_{\perp}, \quad \phi > \phi_{\perp}.
\end{align}
An unrestricted bucket \texttt{ell\_overall} is also saved.

The crucial point is that these buckets use the local basis only to define the
sample subset. The code then computes moments of the \emph{full vector coefficient
magnitude} inside that subset. It does not compute basis-component spectra such as
$P_{\parallel}$, $P_{\xi}$, or $P_{\lambda}$.

\section{Raw coefficient moments and scale-normalized surrogates}

For any coefficient family $W_j$ and any selected bucket $\mathcal{C}_j$, the raw
moment is
\begin{equation}
 M_q^{(j,\mathcal{C})} = \left\langle |W_j|^q \right\rangle_{\mathcal{C}_j}.
\end{equation}
If $W_j$ has units $[X]\,\mathrm{s}^{1/2}$, then
\begin{equation}
 [M_q^{(j,\mathcal{C})}] = [X]^q\,\mathrm{s}^{q/2}.
\end{equation}

The code then forms the saved scale-normalized surrogate
\begin{equation}
 \widetilde{M}_q^{(j,\mathcal{C})} = \frac{M_q^{(j,\mathcal{C})}}{\tau_j^{q/2}}.
\end{equation}
This restores the physical field dimensions $[X]^q$. It does \emph{not} turn the
quantity into a strict finite-difference structure function.

\section{Conditional band-averaged spectra}

For the implemented discretized response, the code evaluates the positive-frequency
response-energy integral on the exact FFT grid actually used in the transform,
\begin{equation}
 A_j^{(\mathrm{disc})} = \sum_{k\ge 0} |\widehat{H}_j(f_k)|^2\,\Delta f.
\end{equation}
In the continuum limit of the implemented normalization,
\begin{equation}
 A_j \to \int_0^{\infty}|\widehat{H}_j(f)|^2\,df = \frac{1}{2\pi},
\end{equation}
but the code uses the discrete form above to match the actual numerical estimator.

The saved second-order spectral observable is
\begin{equation}
 \widehat{P}^{(j,\mathcal{C})} = \frac{\left\langle |W_j|^2 \right\rangle_{\mathcal{C}_j}}{A_j^{(\mathrm{disc})}}.
\end{equation}
Its units are $[X]^2\,\mathrm{s} = [X]^2/\mathrm{Hz}$.

This is a conditional band-averaged PSD estimate associated with the implemented
wavelet response. It is not a Fourier periodogram of a basis-projected component.

\section{Sigma-like quantities and alignment diagnostics}

For two coefficient families $\mathbf{X}_j(t)$ and $\mathbf{Y}_j(t)$, the code
computes the pointwise ratio
\begin{equation}
 \sigma_{\mathrm{ts},j}(t) = \frac{|\mathbf{X}_j(t)|^2 - |\mathbf{Y}_j(t)|^2}{|\mathbf{X}_j(t)|^2 + |\mathbf{Y}_j(t)|^2},
\end{equation}
and the scale-level ratio
\begin{equation}
 \sigma_{\mathrm{scale},j} = \frac{\langle |\mathbf{X}_j|^2 \rangle - \langle |\mathbf{Y}_j|^2 \rangle}{\langle |\mathbf{X}_j|^2 \rangle + \langle |\mathbf{Y}_j|^2 \rangle}.
\end{equation}

For $\mathbf{X}_j=\mathbf{W}_{Z^+,\perp,j}$ and $\mathbf{Y}_j=\mathbf{W}_{Z^-,\perp,j}$,
this is the coefficient-based cross-helicity surrogate. For
$\mathbf{X}_j=\mathbf{W}_{V,\perp,j}$ and $\mathbf{Y}_j=\mathbf{W}_{B_v,\perp,j}$, it is
the coefficient-based residual-energy surrogate.

The physically preferred scale diagnostic is $\sigma_{\mathrm{scale},j}$. The
pointwise quantity is a local surrogate.

\section{Output families and valid interpretation}

The output families should be interpreted as follows.
\begin{itemize}[leftmargin=1.6em,itemsep=0.35em]
\item \texttt{WaveletMoments}: raw coefficient moments with units $[X]^q\,\mathrm{s}^{q/2}$.
\item \texttt{ScaleNormalizedWaveletMoments} (alias of \texttt{Sfuncs}): scale-normalized wavelet moments with units $[X]^q$.
\item \texttt{Spectra} or \texttt{ConditionalBandAveragedPSD}: conditional band-averaged PSD estimates with units $[X]^2/\mathrm{Hz}$.
\item \texttt{B\_nT} and \texttt{B\_vel}: explicit magnetic families in fixed unit systems.
\item \texttt{B}: compatibility family whose unit system still follows \texttt{return\_B\_in\_vel\_units}.
\item \texttt{l\_mag}, \texttt{l\_ell}, \texttt{l\_xi}, and \texttt{l\_lambda}: projected separation magnitudes normalized by $d_i$.
\end{itemize}

\section{What the pipeline does not compute}

The pipeline does not compute exact finite-difference increments. It does not compute
strict two-point structure functions. It does not estimate basis-component spectra by
projecting the fluctuation power onto $\hat{\mathbf{e}}_{\ell}$, $\hat{\mathbf{e}}_{\xi}$,
or $\hat{\mathbf{e}}_{\lambda}$. It does not reconstruct the full 3D spectral tensor
from a single-spacecraft time series.

\section{Bottom-line interpretation}

The revised SCWF pipeline is appropriate when interpreted as a centered,
coefficient-based local-anisotropy analysis built from a Gaussian background and an
even derivative-of-Gaussian band-pass family, with explicit cone-of-influence
masking, consistent lag calibration, conditional band-averaged spectra, and
scale-normalized higher-order wavelet moments.

It becomes physically misleading only when the outputs are overstated as exact
finite-difference structure functions or as basis-component Fourier spectra.

\end{document}
